\scanpage{30}
\mspage{26}

Etant donné un diviseur $D'=\handwritten{\sum}_{P'}m(P').P'$ de $K'$, nous appellerons \uline{extension} de $D'$ à $K$ et noterons $\operatorname{Ext}_{K',K}(D')$ le diviseur de $K$ défini par

\[ \operatorname{Ext}_{K',K}\bigl(\handwritten{\sum}_{P'}m(P').P'\bigr)=\handwritten{\sum}_{P'}m(P').\bigl(\handwritten{\sum}_{P\handwritten{\mid} P'}e(\deleted{\unsure{B|P}}P).P\bigr)\,.\tag*{(31)} \]

Il est clair que $\operatorname{Ext}_{K',K}$ est un \uline{isomorphisme croissant} du groupe \deleted{des} ordonné des diviseurs de $K'$ dans celui des diviseurs de $K$. En vertu de (29) on a

\[ \operatorname{Ext}_{K',K}(D'(x'))=D(x')\quad\text{pour }x'\handwritten{\in} K'\,.\tag*{(32)} \]

Si $K''\handwritten{\subset} K'\handwritten{\subset} K$, et pour tout diviseur $D''$ de $K''$, on a

\[ \operatorname{Ext}_{K',K}(\operatorname{Ext}_{K'',K'}(D''))=\operatorname{Ext}_{K'',K}(D'')\,.\tag*{(33)} \]

Etant donné un diviseur $D=\handwritten{\sum}_P m(P).P$ de $K$, nous appellerons \uline{norme} de $K$ à $K'$ de $D$, et noterons $N_{K/K'}(D)$ le diviseur de $K'$ défini par

\[ N_{K/K'}(D)=\handwritten{\sum}_{P'}\bigl(\handwritten{\sum}_{P\handwritten{\mid} P'}m(P)\underline{f(P)}\bigr).P'\,.\tag*{(34)} \]

Il est clair que $N_{K/K'}$ est un \uline{homomorphisme croissant} du groupe ordonné des diviseurs de $K$ dans celui des diviseurs de $K'$. Comme $\handwritten{\sum}_{P\handwritten{\mid} P'}e(P)f(P)=n$, on voit que

\[ N_{K/K'}(\operatorname{Ext}_{\unsure{K'|\,},K}(D'))=n.D'\,.\tag*{(35)} \]

D'autre part (30) montre que l'on a

\[ D'(N_{K/K'}(x))=N_{K/K'}(D(x))\quad\text{pour tout }x\handwritten{\in} K\,.\tag*{(36)} \]

Enfin, si $K''\handwritten{\subset} K'\handwritten{\subset} K$, et si $D$ est un diviseur de $K$, on a

\[ N_{K'/K''}(N_{K/K'}(D))=N_{K/K''}(D)\,.\tag*{(37)} \]

Supposons maintenant que $\handwritten{\Phi'}$ satisfasse à une formule du produit $\handwritten{\prod}_{P'}(f_{P'}(x'))^{r(P')}=1$. Alors, en vertu de (30), $\handwritten{\Phi}$ satisfait à la formule du produit $\handwritten{\prod}_P(f_P(x))^{r(P)}=1$, où $r(P)=f(P)r(P')$ pour $P\handwritten{\mid} P'$. D'après la définition des \uline{degrés} des diviseurs, et comme $\handwritten{\sum}_{P\handwritten{\mid} P}e(P)f(P)=n$, on a, pour tout couple de diviseurs $D'$ de $K'$ et $D$ de $K$ :

\[ d(\operatorname{Ext}_{K',K}(D'))=n.d(D')\qquad\text{et}\qquad d(N_{K/K'}(D))=d(D)\,.\tag*{(38)} \]

Etant donnée une \uline{répartition} \deleted{$\unsure{X=(x_P)|X=(x_{P'})}$} $X\handwritten{'}=(x'_{P'})$ $(x'_{P'}\handwritten{\in}\hat K'_{P'})$, nous appellerons \uline{extension} de $X'$ à $K$ et noterons $\operatorname{Ext}_{K',K}(X')$ la répartition sur $K$ définie par $X=(x_P)$, où, pour $P\handwritten{\mid} P'$, l'élément $x_P$ de $\hat K_P$ est $x'_{P'}$. Il est clair que l'application $\operatorname{Ext}_{K',K}$ est un \uline{isomorphisme} de l'anneau des répartitions sur $K'$ dans ce
% archive sheet number: 29
